\documentclass[11pt,a4paper]{article}
\usepackage[top=0.75in,bottom=0.75in,left=0.75in,right=0.75in]{geometry}
\usepackage[T1]{fontenc}
\usepackage{times,microtype,enumitem,hyperref,titlesec}
\pagestyle{empty}
\setlength{\parindent}{0pt}
\setlength{\parskip}{0pt}
\linespread{1.10}
\hypersetup{hidelinks,breaklinks=true}
\titleformat{\section}{\bfseries\normalsize}{}{0pt}{}[\titlerule]
\titlespacing*{\section}{0pt}{10pt}{4pt}
\newcommand{\entry}[3]{\noindent\textbf{#1}\hfill\textbf{#3}\\[-1pt]\textit{#2}\par}
\newenvironment{cvitems}{\begin{itemize}[leftmargin=1.1em,label=\textbullet,itemsep=2pt,topsep=3pt,parsep=0pt,partopsep=0pt]}{\end{itemize}}
\begin{document}
{\LARGE\bfseries Aarush Ram Anandh}\par\vspace{4pt}
\href{mailto:aarush.anandh_ug2024@ashoka.edu.in}{aarush.anandh\_ug2024@ashoka.edu.in}\quad|\quad +91 9659786783\par\vspace{2pt}
\href{https://github.com/MushuRam01}{github.com/MushuRam01}\quad|\quad
\href{https://www.linkedin.com/in/aarush-ram-anandh/}{linkedin.com/in/aarush-ram-anandh}\par
\vspace{20pt}
11 November 2026\par\vspace{12pt}
Hiring Team\\
IBM Research / Software Innovation Lab India\par\vspace{12pt}
\textbf{Application: AI Research / AI \& Data Internship}\par\vspace{14pt}
Dear Hiring Team,\par\vspace{10pt}
I am applying for the AI research and AI \& Data tracks of IBM\textquotesingle s 2027 summer internship program in India. I am completing a B.Sc. (Hons.) in Computer Science and Mathematics at Ashoka University, expected in 2028. My experience is in generative models, reinforcement learning and time-series modelling, and I am particularly interested in research that combines controlled experiments with usable implementations.\par\vspace{10pt}
At the Constrained Image-Synthesis Lab, I improved colour fidelity in latent diffusion models to $\Delta E_{00}<1$, with approximately four times better colour accuracy than evaluated open-source baselines. This work is reported in the public preprint \textit{Inference-Time Loss-Guided Colour Preservation in Diffusion Sampling} (arXiv:2601.17259). I am also developing Controlled Generalised Sliced Wasserstein, a controllable optimal-transport framework with a paper submitted to JMLR; my work includes computational methods in C++ and specifications and proofs in Lean.\par\vspace{10pt}
In a two-person RL coursework project, we co-developed a Gymnasium/poke-env environment with legal-action masks, potential-based reward shaping, MaskablePPO and masked DQN agents. We evaluated PPO over two variants and five seeds each, reaching 99.7--100\% wins against the fixed heuristic over 1,020 final battles per run. In a separate team project with Tanush Seal, I worked on financial forecasting using a CatBoost baseline and HMM, Transformer and other residual models, with company-level diagnostic dashboards. At the Koita Centre for Digital Health, I contributed to a wearable-data pipeline achieving approximately 70\% sensitivity for asthma exacerbation detection in study evaluation, and built experimental frozen time-series foundation-model routes for surprise features and embeddings. Those foundation-model routes were implemented but not yet run in the reported cohort evaluation.\par\vspace{10pt}
At Avartan Labs, I raised rare-class correct detection from 40\% to 80\% in an RT-DETR pipeline covering 18 material classes. This industry internship gave me experience analysing detector failures and investigating targeted model improvements.\par\vspace{10pt}
The emphasis on model behaviour, reinforcement learning, data and experimental research in IBM\textquotesingle s AI tracks fits the direction I want to pursue. I would welcome the opportunity to contribute Python/PyTorch implementation and evaluation work, while learning how these methods support IBM\textquotesingle s research and enterprise applications. Thank you for considering my application.\par\vspace{14pt}
Yours sincerely,\\[5pt]
Aarush Ram Anandh
\end{document}
